\begin{frame}{$P$ 추정 결과: 모델 기반 vs 모델 없는}
  2000에피소드의 무작위 탐색으로 추정된 $P$ 는, 실제 FrozenLake
  구현(\texttt{is\_slippery=True}일 때 intended 방향 2/3,
  좌/우 1/6씩)에 \textbf{수렴}.
  \vskip0.6em
  \begin{block}{핵심: ``$P$ 의 사용 방식''이 알고리즘의 가족을
  나눈다}
    이 $P_{\text{est}}$ 를 Week 4의 동적계획법에 그대로 대입하면
    \kb{모델 기반} 알고리즘을 돌릴 수 있고, $P$ 를 무시하고
    \textbf{경험만} 쓰면 \kb{모델 없는} 알고리즘(Week 6--)을
    돌릴 수 있다.
    \vskip0.3em
    {\small \kq{같은 환경, 같은 샘플, 다른 ``$P$ 의 사용 방식''이
    알고리즘의 가족을 나눈다.}}
  \end{block}
  \vskip0.5em
  {\footnotesize 모델 기반(Week 4--5): 추정된 $P$ 를 명시적으로
  써서 계획. 모델 없는(Week 6--11): $P$ 를 전혀 쓰지 않고
  경험(실제 전이 샘플)에서만 가치를 배운다.
  이 ``첫 갈림길''이 이 학기의 알고리즘 지도를 결정한다.}
\end{frame}
